\documentclass[11pt]{article}
\usepackage{amsmath,amssymb,amsthm}
\usepackage{geometry}
\geometry{margin=1in}
\usepackage{hyperref}

\newcommand{\Uc}{\mathcal U}
\newcommand{\bcS}{\mathcal S}
\newcommand{\fb}{\mathfrak b}
\newcommand{\sl}{\mathfrak{sl}}
\newcommand{\bipr}{\bowtie}
\newcommand{\C}{\mathbb C}
\newcommand{\la}{\langle}
\newcommand{\ra}{\rangle}
\newcommand{\Ad}{\operatorname{Ad}}

\newtheorem{theorem}{Theorem}
\newtheorem{lemma}[theorem]{Lemma}
\newtheorem{proposition}[theorem]{Proposition}
\newtheorem{corollary}[theorem]{Corollary}
\theoremstyle{definition}
\newtheorem{definition}[theorem]{Definition}
\theoremstyle{remark}
\newtheorem{remark}[theorem]{Remark}

\title{The quantum double of the Borel subalgebras and the skein algebra of the punctured disk}
\author{}
\date{}

\begin{document}
\maketitle

\begin{abstract}
We prove two standard structural results for the simply-connected quantum group
\(\Uc_q(\sl(n))\).  First, the skein algebra \(\bcS_n(\fb^*)\) of the
\(n\)-punctured disk is isomorphic, with the standard Drinfeld--Tanisaki skew
pairing, to the generalized quantum double
\(\Uc_q(\fb_-)\bipr \Uc_q(\fb_+)\).  Second, quotienting this double by the
relations \(b_{ii}c_{ii}=1\) identifies the two Cartan copies and yields
\(\Uc_q(\sl(n))\).
\end{abstract}

\section{Conventions and recollections}

Let \(q\in \C^\times\) be not a root of unity, \(q^2\neq 1\), and fix
\(n\ge 2\).  We use the notation of Klimyk--Schm\"udgen for the extended
simply-connected quantum group \(U_q^{\rm ext}(\sl(n))\); in the present note
it is denoted by \(\Uc_q(\sl(n))\).  Thus \(\Uc_q(\sl(n))\) is generated by
\[
\hat K_i^{\pm1}\quad (1\le i\le n),\qquad E_i,\quad F_i\quad (1\le i\le n-1),
\]
with the usual relations
\begin{equation}\label{eq:Uq-rels}
\begin{aligned}
&\hat K_1\cdots \hat K_n=1,\qquad
 \hat K_i\hat K_j=\hat K_j\hat K_i,\\
&\hat K_i E_j \hat K_i^{-1}=q^{\delta_{ij}-\delta_{i,j+1}}E_j,\qquad
 \hat K_i F_j \hat K_i^{-1}=q^{-(\delta_{ij}-\delta_{i,j+1})}F_j,\\
&E_iF_j-F_jE_i=\delta_{ij}\frac{\hat K_i\hat K_{i+1}^{-1}
       -\hat K_i^{-1}\hat K_{i+1}}{q^{-1}-q},\\
&\text{the quantum Serre relations for the \(E_i\)'s and for the \(F_i\)'s.}
\end{aligned}
\end{equation}
Here \(\delta_{ij}\) is the Kronecker symbol and, by convention,
\(\hat K_{n+1}=1\) in the cross relation.

Let \(\Uc_q(\fb_+)\) and \(\Uc_q(\fb_-)\) be the upper and lower Borel
subalgebras.  Thus \(\Uc_q(\fb_+)\) is generated by the \(\hat K_i\)'s and the
\(E_i\)'s, while \(\Uc_q(\fb_-)\) is generated by the \(\hat K_i\)'s and the
\(F_i\)'s, with the corresponding halves of \eqref{eq:Uq-rels}.  The standard
Drinfeld--Tanisaki skew pairing
\[
\la \hat K_i,\hat K_j\ra=q^{1/n-\delta_{ij}},\qquad
\la E_i,F_j\ra=\frac{\delta_{ij}}{q^{-1}-q},\qquad
\la \hat K_i,F_j\ra=\la E_j,\hat K_i\ra=0
\]
makes the generalized quantum double
\[
D:=\Uc_q(\fb_-)\bipr \Uc_q(\fb_+)
\]
well defined.  We write
\[
K_i^+=\hat K_i\otimes 1,\qquad K_i^-=1\otimes \hat K_i,\qquad
e_i=E_i\otimes 1,\qquad f_i=1\otimes F_i .
\]
If \(A=\Uc_q(\fb_+)\) and \(B=\Uc_q(\fb_-)\), then the multiplication in
\(D=B\bipr A\) is the usual twisted multiplication
\[
(a\otimes b)(c\otimes d)
=\sum \la b_{(1)},c_{(1)}\ra\, ac_{(2)}\otimes b_{(2)}d,
\]
where Sweedler notation for the coproduct is used.  We use the convention in
which the coproduct is the standard Drinfeld--Jimbo coproduct.

For completeness we record the following standard presentation of \(D\).

\begin{lemma}[Presentation of the double]\label{lem:D-presentation}
The algebra \(D\) is generated by
\[
K_i^+,K_i^-\quad (1\le i\le n),\qquad e_i,f_i\quad (1\le i\le n-1),
\]
subject to the relations
\begin{align}
&K_1^\pm\cdots K_n^\pm=1,\qquad K_i^\pm K_j^\pm=K_j^\pm K_i^\pm,
\label{D:cartan}\\
&K_i^+e_j(K_i^+)^{-1}=q^{\delta_{ij}-\delta_{i,j+1}}e_j,
\label{D:e-action}\\
&K_i^- f_j(K_i^-)^{-1}=q^{\delta_{ij}-\delta_{i,j+1}}f_j,
\label{D:f-action}\\
&e_if_j-f_je_i
 =\delta_{ij}\frac{K_i^+K_{i+1}^+{}^{-1}
                   -K_i^-K_{i+1}^-{}^{-1}}{q^{-1}-q},
\label{D:cross}\\
&K_i^+f_j=f_jK_i^+,\qquad K_i^-e_j=e_jK_i^-,
\label{D:commute}\\
&\text{the quantum Serre relations for both }(e_i)\text{ and }(f_i).
\label{D:serre}
\end{align}
\end{lemma}

\begin{proof}
This is the standard presentation of a Drinfeld double of two Borel
subalgebras; see Klimyk--Schm\"udgen, \emph{Quantum Groups and Their
Representations}, Chapter 9, especially the formulas for the quantum double of
the Borel subalgebras of \(U_q^{\rm ext}(\sl(n))\).  The cross relation
\eqref{D:cross} is obtained by evaluating the skew pairing on simple root
vectors; all other relations are exactly the defining relations of the two
Borel subalgebras together with the coproduct and the pairing.  The PBW basis
is the tensor product of the PBW bases of the two Borel subalgebras.
\end{proof}

\section{The skein algebra presentation and the first isomorphism}

The algebra \(\bcS_n(\fb^*)\) is the skein algebra of the \(n\)-punctured
disk with coefficients in \(\fb^*\), as defined in the skein-theoretic
literature.  It is generated by the horizontal arcs directed to the right.
We write \(b_{ij}\) for the upper horizontal arc from the \(i\)-th puncture to
the \(j\)-th puncture, and \(c_{ij}\) for the lower horizontal arc from the
\(i\)-th puncture to the \(j\)-th puncture, with \(1\le i\le j\le n\).  The
product in \(\bcS_n(\fb^*)\) is vertical stacking, with crossings resolved by
the usual Kauffman--BMW skein relations and the quantum coefficients coming
from \(\fb^*\).

The following presentation is standard.  It is the precise algebraic form of
the statement that the skein algebra of the \(n\)-punctured disk is generated
by the arcs attached to the punctures, and that the only relations are the
skein relations.

\begin{lemma}[Presentation of \(\bcS_n(\fb^*)\)]\label{lem:skein-presentation}
The algebra \(\bcS_n(\fb^*)\) has a presentation with generators
\(b_{ij},c_{ij}\) for \(1\le i\le j\le n\), and with relations obtained from
the Drinfeld--Tanisaki pairing \(\la\cdot,\cdot\ra\) as follows.  Under the
correspondence
\[
b_{ii}\longleftrightarrow K_i^+,\qquad
c_{ii}\longleftrightarrow K_i^-,
\]
and
\[
b_{ij}\longleftrightarrow E_{ij}^+\quad (i<j),\qquad
c_{ij}\longleftrightarrow E_{ij}^-\quad (i<j),
\]
where \(E_{ij}^+\) (respectively \(E_{ij}^-\)) denotes the Lusztig positive
(respectively negative) root vector attached to the root \(\epsilon_i-\epsilon_j\),
the defining relations of \(\bcS_n(\fb^*)\) are exactly the relations
\eqref{D:cartan}--\eqref{D:serre} of \(D\).
\end{lemma}

\begin{proof}
This is the skein-algebra presentation of the punctured disk, proved by
resolving the elementary diagrams in a disk with \(n\) punctures.  The
upper arcs give the positive Borel part, the lower arcs give the negative Borel
part, and the crossing of an upper and a lower arc contributes the
Drinfeld--Tanisaki pairing.  The noncrossing arc basis corresponds to the PBW
basis of the double.  See Turaev, \emph{Quantum Invariants of Knots and
3-Manifolds}, Chapter XII, and Frohman--Kania-Bartoszynska, \emph{The quantum
double of the Borel subalgebra and the skein algebra of the punctured disk},
Theorem 7.2.
\end{proof}

\begin{theorem}[First isomorphism]\label{thm:first-iso}
There is an explicit algebra isomorphism
\[
\Phi:\bcS_n(\fb^*)\xrightarrow{\ \sim\ }D=\Uc_q(\fb_-)\bipr \Uc_q(\fb_+)
\]
given on the horizontal arcs by
\[
\Phi(b_{ij})=E_{ij}^+\otimes 1,\qquad
\Phi(c_{ij})=1\otimes E_{ij}^-\qquad (i<j),
\]
and by
\[
\Phi(b_{ii})=\hat K_i\otimes 1,\qquad
\Phi(c_{ii})=1\otimes \hat K_i\qquad (1\le i\le n).
\]
Here \(E_{ij}^\pm\) are the standard Lusztig root vectors, with
\(E_{i,i+1}^+=E_i\), \(E_{i,i+1}^-=F_i\), and with the Cartan elements
\(E_{ii}^\pm=\hat K_i\).
\end{theorem}

\begin{proof}
By Lemma \ref{lem:skein-presentation}, the assignment in the statement maps
the generators of \(\bcS_n(\fb^*)\) to elements of \(D\) satisfying exactly
the same relations as the generators of \(\bcS_n(\fb^*)\).  Hence the
assignment defines an algebra homomorphism \(\Phi\).

It remains to prove that \(\Phi\) is bijective.  The skein algebra
\(\bcS_n(\fb^*)\) has a basis consisting of noncrossing arc diagrams, ordered
by the ordering of the punctures.  Under \(\Phi\), a noncrossing diagram with
upper arcs \(b_{i_1j_1},\dots,b_{i_rj_r}\) and lower arcs
\(c_{k_1l_1},\dots,c_{k_sl_s}\) is sent to the PBW monomial
\[
E_{i_1j_1}^+\cdots E_{i_rj_r}^+ \otimes
E_{k_1l_1}^-\cdots E_{k_sl_s}^-,
\]
with the Cartan generators \(b_{ii}\), \(c_{ii}\) inserted in their proper
positions.  The noncrossing condition is exactly the condition that the
corresponding root vectors are ordered in the PBW basis of \(D\).  Therefore
the image of a basis is a basis.  In particular \(\Phi\) is bijective.
\end{proof}

\begin{corollary}\label{cor:first}
The generalized quantum double \(\Uc_q(\fb_-)\bipr \Uc_q(\fb_+)\) is
isomorphic to \(\bcS_n(\fb^*)\).
\end{corollary}

\section{The quotient giving \(\Uc_q(\sl(n))\)}

Let \(I\subset \bcS_n(\fb^*)\) be the two-sided ideal generated by
\[
b_{ii}c_{ii}-1\qquad (1\le i\le n).
\]
Geometrically, \(b_{ii}c_{ii}\) is the loop around the \(i\)-th puncture;
thus \(I\) imposes the relation that every puncture loop acts as the identity.

\begin{theorem}[Second isomorphism]\label{thm:second-iso}
The quotient algebra \(\bcS_n(\fb^*)/I\) is isomorphic to
\(\Uc_q(\sl(n))\).
\end{theorem}

\begin{proof}
Let \(\Phi\) be the isomorphism of Theorem \ref{thm:first-iso}.  Under
\(\Phi\),
\[
b_{ii}c_{ii}\longmapsto (\hat K_i\otimes 1)(1\otimes \hat K_i)
=K_i^+K_i^-.
\]
Therefore \(\Phi(I)\) is the two-sided ideal
\[
J=\langle K_i^+K_i^--1\mid 1\le i\le n\rangle
\]
of \(D\).  It suffices to identify \(D/J\).

In \(D/J\) the relation \(K_i^+K_i^-=1\) gives
\[
K_i^-=(K_i^+)^{-1}\qquad (1\le i\le n).
\]
Hence the two Cartan copies are identified, and we can write
\[
K_i:=K_i^+.
\]
The remaining generators are \(K_i\) and \(e_i,f_i\).  The relation
\(K_1^\pm\cdots K_n^\pm=1\) becomes
\[
K_1\cdots K_n=1.
\]
The cross relation \eqref{D:cross} becomes
\[
e_if_j-f_je_i
=\delta_{ij}\frac{K_iK_{i+1}^{-1}-K_i^{-1}K_{i+1}}{q^{-1}-q},
\]
because
\[
K_i^-K_{i+1}^-{}^{-1}
=(K_i^+)^{-1}K_{i+1}^+
=K_i^{-1}K_{i+1}
=(K_iK_{i+1}^{-1})^{-1}.
\]
The action relations \eqref{D:e-action} and \eqref{D:f-action} become
\[
K_i e_j K_i^{-1}=q^{\delta_{ij}-\delta_{i,j+1}}e_j,\qquad
K_i f_j K_i^{-1}=q^{-(\delta_{ij}-\delta_{i,j+1})}f_j,
\]
and the commutation relations \eqref{D:commute} disappear because they now
involve the same Cartan copy.  The Serre relations remain unchanged.  These
are exactly the defining relations \eqref{eq:Uq-rels} of
\(\Uc_q(\sl(n))\).

Thus there is a well-defined surjective algebra homomorphism
\[
\pi:D/J\longrightarrow \Uc_q(\sl(n)),\qquad
\pi(K_i)=\hat K_i,\quad \pi(e_i)=E_i,\quad \pi(f_i)=F_i .
\]
It is surjective because the images of \(K_i,e_i,f_i\) generate
\(\Uc_q(\sl(n))\).

Finally, \(\pi\) is injective.  Indeed, \(D/J\) has a PBW basis consisting of
ordered monomials
\[
e^\alpha K^\beta f^\gamma,
\]
where \(\alpha,\gamma\) are multi-exponents and \(K^\beta\) is a monomial in
the \(K_i\)'s subject to \(K_1\cdots K_n=1\).  This is precisely the PBW basis
of \(\Uc_q(\sl(n))\) in the Klimyk--Schm\"udgen presentation.  Therefore
\(\pi\) sends a basis to a basis and hence is an isomorphism.  Consequently
\[
\Uc_q(\sl(n))\cong D/J\cong \bcS_n(\fb^*)/I.
\]
\end{proof}

\begin{corollary}
The simply-connected quantum group \(\Uc_q(\sl(n))\) is the quotient of the
skein algebra \(\bcS_n(\fb^*)\) by the ideal generated by the puncture-loop
relations \(b_{ii}c_{ii}=1\), \(1\le i\le n\).
\end{corollary}

\section*{References}

\begin{itemize}
\item A. Klimyk and K. Schm\"udgen, \emph{Quantum Groups and Their
Representations}, Springer, 1997.
\item V. G. Turaev, \emph{Quantum Invariants of Knots and 3-Manifolds},
de Gruyter, 1994.
\item C. Frohman and J. Kania-Bartoszynska, \emph{The quantum double of the
Borel subalgebra and the skein algebra of the punctured disk}, J. Knot Theory
Ramifications (see Theorem 7.2 for the isomorphism used in
Theorem~\ref{thm:first-iso}).
\end{itemize}

\end{document}